%==============================================================================
% Chapitre 38 - Stabilité dynamique
%==============================================================================

\section{Introduction}

La stabilité gyroscopique et la stabilité dynamique sont souvent confondues.

Pourtant, elles décrivent deux phénomènes différents.

La stabilité gyroscopique répond essentiellement à la question :

\begin{quote}
Le projectile tourne-t-il suffisamment vite pour éviter le basculement ?
\end{quote}

La stabilité dynamique s'intéresse quant à elle à l'évolution du mouvement
du projectile au cours du temps.

\begin{aretenir}

Un projectile peut être gyroscopiquement stable tout en étant
dynamiquement instable.

\end{aretenir}

\section{Rappel sur la stabilité gyroscopique}

La stabilité gyroscopique dépend principalement :

\begin{itemize}
    \item du pas de rayure ;
    \item de la vitesse de rotation ;
    \item de la géométrie ;
    \item des caractéristiques aérodynamiques.
\end{itemize}

Elle est souvent caractérisée par le facteur :

\begin{equation}
S_g
\end{equation}

que nous avons étudié précédemment.

\section{Pourquoi cela ne suffit-il pas ?}

Pendant son vol, un projectile est soumis à :

\begin{itemize}
    \item des perturbations aérodynamiques ;
    \item des variations de vitesse ;
    \item des changements de régime de vol ;
    \item des mouvements oscillatoires.
\end{itemize}

Ces phénomènes évoluent continuellement.

La seule connaissance de $S_g$ ne permet donc pas de décrire complètement
le comportement du projectile.

\section{Les oscillations du projectile}

Nous avons déjà rencontré :

\begin{itemize}
    \item la précession ;
    \item la nutation.
\end{itemize}

Ces mouvements produisent des oscillations de l'axe du projectile.

La stabilité dynamique décrit l'évolution de ces oscillations.

\section{Oscillations amorties}

Dans un projectile dynamiquement stable :

\begin{itemize}
    \item les oscillations diminuent progressivement ;
    \item l'axe se stabilise ;
    \item le comportement devient plus régulier.
\end{itemize}

On parle d'amortissement.

\section{Oscillations entretenues}

Dans certains cas :

\begin{itemize}
    \item les oscillations ne diminuent pas ;
    \item leur amplitude reste importante ;
    \item la précision se dégrade.
\end{itemize}

Le comportement devient alors moins favorable.

\section{Oscillations divergentes}

Dans les cas les plus défavorables :

\begin{itemize}
    \item les oscillations augmentent ;
    \item les angles d'incidence deviennent importants ;
    \item le projectile peut perdre sa stabilité.
\end{itemize}

Cette situation correspond à une instabilité dynamique.

\section{Analogie avec un amortisseur}

Un véhicule constitue un exemple simple.

Après une bosse :

\begin{itemize}
    \item un bon amortisseur réduit rapidement les oscillations ;
    \item un mauvais amortisseur les laisse persister.
\end{itemize}

La stabilité dynamique joue un rôle similaire pour les projectiles.

\section{Influence de la vitesse}

La stabilité dynamique dépend fortement de la vitesse.

Au cours du vol :

\begin{itemize}
    \item la vitesse diminue ;
    \item les forces aérodynamiques évoluent ;
    \item les coefficients aérodynamiques changent.
\end{itemize}

Le comportement dynamique peut donc se modifier.

\section{Influence du nombre de Mach}

Le nombre de Mach joue un rôle important.

Les propriétés aérodynamiques diffèrent fortement selon que le projectile est :

\begin{itemize}
    \item supersonique ;
    \item transsonique ;
    \item subsonique.
\end{itemize}

Certaines instabilités apparaissent principalement dans la zone
transsonique.

\section{Passage transsonique}

Le passage transsonique constitue souvent une phase critique.

Pendant cette transition :

\begin{itemize}
    \item les ondes de choc évoluent ;
    \item les centres aérodynamiques se déplacent ;
    \item les coefficients aérodynamiques varient rapidement.
\end{itemize}

Le comportement du projectile peut être affecté.

\begin{aretenir}

Le domaine transsonique constitue souvent la zone la plus délicate du vol.

\end{aretenir}

\section{Influence de la géométrie}

La stabilité dynamique dépend également :

\begin{itemize}
    \item de l'ogive ;
    \item du boat-tail ;
    \item de la longueur ;
    \item de la répartition des masses.
\end{itemize}

Deux projectiles possédant le même facteur de stabilité gyroscopique peuvent
présenter des comportements dynamiques très différents.

\section{Influence du centre de gravité}

La position du centre de gravité influence directement :

\begin{itemize}
    \item les moments aérodynamiques ;
    \item la précession ;
    \item la nutation ;
    \item l'amortissement.
\end{itemize}

Elle constitue un paramètre majeur de conception.

\section{Influence du centre de pression}

Le centre de pression intervient également.

Son déplacement au cours du vol peut modifier les caractéristiques
dynamiques du projectile.

Cette évolution est particulièrement importante à haute vitesse.

\section{Conséquences pratiques}

Une mauvaise stabilité dynamique peut entraîner :

\begin{itemize}
    \item une augmentation de la dispersion ;
    \item une perte de précision ;
    \item une augmentation de la traînée ;
    \item des écarts imprévus.
\end{itemize}

Ces phénomènes deviennent particulièrement visibles aux longues distances.

\section{Mesure expérimentale}

La stabilité dynamique est généralement étudiée à l'aide :

\begin{itemize}
    \item de radars Doppler ;
    \item de caméras ultra-rapides ;
    \item d'essais instrumentés ;
    \item de simulations numériques.
\end{itemize}

Son évaluation est plus complexe que celle de la stabilité gyroscopique.

\section{Stabilité dynamique et simulation}

Les modèles simples représentent difficilement ces phénomènes.

Les modèles avancés permettent :

\begin{itemize}
    \item de calculer les oscillations ;
    \item d'estimer leur amortissement ;
    \item d'évaluer les risques d'instabilité.
\end{itemize}

Les modèles 6DOF sont particulièrement adaptés à cette analyse.

\begin{simulation}

L'étude de la stabilité dynamique constitue l'une des principales raisons
justifiant l'utilisation des modèles 6DOF.

\end{simulation}

\section{Vers les critères de stabilité modernes}

Nous avons maintenant étudié :

\begin{itemize}
    \item la stabilité statique ;
    \item la stabilité gyroscopique ;
    \item la stabilité dynamique.
\end{itemize}

Ces notions permettent de comprendre pourquoi certains projectiles offrent
d'excellentes performances alors que d'autres deviennent rapidement
imprécis.

Les chapitres suivants aborderont des critères et méthodes plus avancés
utilisés dans la conception moderne des projectiles.

\section{À retenir}

\begin{aretenir}

\begin{itemize}
    \item La stabilité dynamique décrit l'évolution des oscillations dans le
          temps.
    \item Un projectile peut être gyroscopiquement stable mais
          dynamiquement instable.
    \item Les phénomènes transsoniques jouent souvent un rôle important.
    \item La géométrie et la répartition des masses influencent fortement la
          stabilité dynamique.
    \item Les modèles avancés sont souvent nécessaires pour l'étudier.
\end{itemize}

\end{aretenir}
